\section{Compiler Passes and Memory Plan}
\label{app:compiler}

\Cref{sec:method} presents the compiled representation by what it contains;
this appendix presents it by how it is produced. \Cref{alg:compile} lists the
logical passes and \cref{fig:compiler} draws them on the example game of
\cref{fig:archpipeline}; these are data dependencies.
Node and edge walks plus the stable depth sort cost
$O(\numnodes + \numedges \log \numedges)$ time and $O(\numnodes + \numedges)$
host storage; device allocations may precede planning. \Cref{tab:overhead} and
\cref{tab:lifecycle} report the measured build times; on the 434k-state turn
subgame the game specification dominates and compilation proper is a fraction
of a second.

\begin{algorithm}[t]
\caption{\textsc{Compile}(GameSpec $G$, rule $\rho$, device $d$) $\to$ solver}
\label{alg:compile}
\begin{algorithmic}[1]
\State \textbf{Validate} $G$: node ids contiguous, one parent per non-root node
(tree), every information set owned by one player and its action list equal to
each member node's, chance probabilities present for every chance edge.
\Comment{structural checks}
\State \textbf{Slots}: sort information sets by key; $\mathrm{off}(I) \gets$
prefix sum of $|\acts(I)|$; $\numslots \gets \sum_I |\acts(I)|$.
\State \textbf{Node arrays}: number nodes parent-before-child; emit type, owner,
information set, action offset/count, terminal payoff $u$, chance probability
$p_c$ of each chance action.
\State \textbf{Edge arrays}: for each edge $(h,a,h')$ emit parent $h$, child
$h'$, slot $\mathrm{off}(I(h)) + a$ or the sentinel for chance edges, owner
flags, sign $s(e)$, multiplier $p_c$ or $1$. \Comment{\cref{sec:compiled}}
\State \textbf{Depth schedule}: $\mathrm{depth}(h)$ by one forward walk; stable
sort edges by $\mathrm{depth}(\text{parent})$; block $b$ = all edges of one
depth; assert parent depth $<$ child depth for every edge.
\Comment{$D$ blocks, \cref{sec:dataflow}(b)}
\State \textbf{Chance folding}: walk the blocks once, $\pi_c(h') \gets
\pi_c(h)\, p_c(h,a,h')$; template $\cfv_{\mathrm{tmpl}}(z) \gets \pi_c(z)\,u(z)$;
retain the root action owning each node and $\pi_c$ with the root factor
removed, for later range updates. \Comment{\cref{eq:template}}
\State \textbf{Index emission}: per block, dual-lane $P_2, C_2, S_2$ (chance and
other-player entries point at the sentinel) and single-lane $P_{\mathrm{val}},
C_{\mathrm{val}}, S_{\mathrm{val}}$; flat over decision edges $P_{\mathrm{reg}},
C_{\mathrm{reg}}, S_{\mathrm{reg}}, O_{\mathrm{reg}}, M_{\mathrm{reg}}, s$; per
player, the halves of the flat arrays for alternating updates.
\State \textbf{Memory plan}: estimate scheduled storage from
$\numnodes, \numslots, \numinfosets$ and dtype; record budget fit as a
planning indicator, provision iteration buffers, set root reach to one.
\State \textbf{Rule lowering}: bind $\rho$'s hooks into the iteration body:
\textsc{discount}$(\regret,\bar{s},t)$ before the passes,
\textsc{after}$(\regret)$ after them, and a prediction buffer when $\rho$
matches against $\regret + r_t$. \Comment{\cref{tab:rules}}
\State \textbf{Runtime}: on CUDA, three eager warmup iterations, capture of one
(or $k$) iteration bodies, record of the persistent buffer addresses; replay
thereafter with the device-side counter refilled per call.
\Comment{\cref{sec:cudagraph}}
\end{algorithmic}
\end{algorithm}

% Arch-3 (appendix): the compiler passes of Alg. 2 drawn on the same 11-node
% example game as Figures 2 and 3, with the same glyph and color vocabulary
% (gray chance, teal P1, violet P2, dark terminal / sentinel, orange = commit).
% Full text width. Top lane: host-side passes 1-7 turn the tree into flat
% strips, sort the edges into depth blocks, fold chance into the template, and
% emit per-lane index arrays. Bottom lane: passes 8-10 size and pin the buffers,
% bind the rule hooks into the iteration body, warm up, record, and replay; the
% dashed orange return path shows that between solves only mutable buffers are
% rewritten, so the captured graph survives. Badges are Alg. 2 line numbers.
\begin{figure*}[!htb]
\centering
\resizebox{\textwidth}{!}{%
\begin{tikzpicture}[
    font=\sffamily\small,
    dec1/.style={circle, draw=teal!60!black, fill=teal!45, minimum size=2.6mm, inner sep=0pt},
    dec2/.style={circle, draw=violet!60!black, fill=violet!40, minimum size=2.6mm, inner sep=0pt},
    chn/.style={diamond, draw=black!65, fill=black!22, minimum size=3.2mm, inner sep=0pt},
    ter/.style={rectangle, draw=black!70, fill=black!55, minimum size=2.0mm, inner sep=0pt},
    edge/.style={black!55, line width=0.7pt},
    flow/.style={-{Stealth[length=2.2mm]}, line width=1.0pt, black!60},
    arr/.style={-{Stealth[length=1.6mm]}, line width=0.9pt, black!60},
    fold/.style={-{Stealth[length=1.8mm]}, densely dashed, line width=0.8pt, black!55},
    ret/.style={-{Stealth[length=1.8mm]}, densely dashed, line width=0.9pt, orange!75!black},
    hook/.style={draw=orange!75!black, densely dashed, rounded corners=1.5pt, inner sep=1.5pt, text=orange!60!black},
    tick/.style={orange!80!black, line width=1.0pt},
    gtick/.style={orange!95!black, line width=1.6pt},
    note/.style={font=\sffamily\small, text=black!55},
    lanelabel/.style={font=\sffamily\small\bfseries, text=black!60},
    hdr/.style={font=\sffamily\small, text=black!65, anchor=west},
    badge/.style={circle, fill=orange!85!black, text=white,
                  font=\sffamily\small\bfseries, inner sep=1.0pt}
]
\def\cw{0.34}
% one strip cell
\newcommand{\cell}[3]{\draw[black!60, fill=#3] ({#1},{#2}) rectangle ++(\cw,\cw);}
% separators at cell boundaries: x, y-bottom, list
\newcommand{\seps}[3]{\foreach \s in {#3}
  \draw[black!85, line width=1.0pt] ({#1+\s*\cw},#2-0.03) -- ++(0,\cw+0.06);}
% the 10 edges in a depth-block layout (2 | 4 | 4) with gaps: x, y-bottom, fills
\newcommand{\blockstrip}[3]{%
  \foreach [count=\i from 0] \c in #3{
    \pgfmathsetmacro{\g}{(\i>1)*0.12 + (\i>5)*0.12}
    \cell{#1+\i*\cw+\g}{#2}{\c}}}
% owner fill lists for the 10 edges (BFS order: e0,e1 chance; e2-e5 P1; e6-e9 P2)
\def\Pfills{black!22, black!22, teal!45, teal!45, teal!45, teal!45, violet!40, violet!40, violet!40, violet!40}
\def\Cfills{teal!45, teal!45, violet!40, black!55, violet!40, black!55, black!55, black!55, black!55, black!55}
\def\Sfills{black!45, black!45, teal!45, teal!45, teal!45, teal!45, violet!40, violet!40, violet!40, violet!40}
\def\Qone{black!45, black!45, teal!45, teal!45, teal!45, teal!45, black!45, black!45, black!45, black!45}
\def\Qtwo{black!45, black!45, black!45, black!45, black!45, black!45, violet!40, violet!40, violet!40, violet!40}
\def\Nfills{black!22, teal!45, teal!45, violet!40, black!55, violet!40, black!55, black!55, black!55, black!55, black!55}
\def\Tfills{white, white, white, white, black!55, white, black!55, black!55, black!55, black!55, black!55}

% ======================= host lane: passes 1-7 =======================
% --- the example game, input to the compiler ---
\node[chn]  (c0)  at (1.45,-0.70) {};
\node[dec1] (p1a) at (0.85,-1.40) {};
\node[dec1] (p1b) at (2.05,-1.40) {};
\node[dec2] (q2a) at (0.50,-2.10) {};
\node[ter]  (t2a) at (1.20,-2.10) {};
\node[dec2] (q2b) at (1.70,-2.10) {};
\node[ter]  (t2b) at (2.40,-2.10) {};
\node[ter]  (t3a) at (0.30,-2.80) {};
\node[ter]  (t3b) at (0.70,-2.80) {};
\node[ter]  (t3c) at (1.50,-2.80) {};
\node[ter]  (t3d) at (1.90,-2.80) {};
\draw[edge, densely dashed] (c0) -- (p1a); \draw[edge, densely dashed] (c0) -- (p1b);
\draw[edge] (p1a) -- (q2a); \draw[edge] (p1a) -- (t2a);
\draw[edge] (p1b) -- (q2b); \draw[edge] (p1b) -- (t2b);
\draw[edge] (q2a) -- (t3a); \draw[edge] (q2a) -- (t3b);
\draw[edge] (q2b) -- (t3c); \draw[edge] (q2b) -- (t3d);
\node[note] at (1.45,-3.35) {game specification $G$};
% --- pass 1: validation gate ---
\node[badge] at (3.05,-1.05) {1};
\node[diamond, draw=black!65, fill=blue!6, minimum size=5.2mm, inner sep=0pt] (gate) at (3.05,-1.65) {};
\draw[green!45!black, line width=1.0pt] (2.93,-1.65) -- (3.02,-1.76) -- (3.18,-1.54);
\draw[arr] (2.55,-1.65) -- (gate.west);
\draw[arr] (gate.east) -- (3.60,-1.65);
% --- passes 2-4: flat strips ---
\def\sxa{3.65}
\node[badge] at (3.78,-0.42) {2};
\node[hdr] at (3.98,-0.42) {slots: prefix sum over sets};
\foreach [count=\i from 0] \c in {teal!45, teal!45, teal!45, teal!45, violet!40, violet!40, violet!40, violet!40}
  \cell{\sxa+\i*\cw}{-0.95}{\c}
\seps{\sxa}{-0.95}{0, 2, 4, 6, 8}
\node[badge] at (3.78,-1.45) {3};
\node[hdr] at (3.98,-1.45) {nodes, parent before child};
\foreach [count=\i from 0] \c in \Nfills \cell{\sxa+\i*\cw}{-1.98}{\c}
\seps{\sxa}{-1.98}{0, 1, 3, 7, 11}
\node[badge] at (3.78,-2.48) {4};
\node[hdr] at (3.98,-2.48) {edges: parent, child, slot};
\foreach [count=\i from 0] \c in \Pfills \cell{\sxa+\i*\cw}{-3.00}{\c}
\foreach [count=\i from 0] \c in \Cfills \cell{\sxa+\i*\cw}{-3.34}{\c}
\foreach [count=\i from 0] \c in \Sfills \cell{\sxa+\i*\cw}{-3.68}{\c}
\node[note, anchor=east] at (3.57,-2.83) {$P$};
\node[note, anchor=east] at (3.57,-3.17) {$C$};
\node[note, anchor=east] at (3.57,-3.51) {$S$};
% --- pass 6 (drawn above 5): chance folding into the values template ---
\def\sxb{7.95}
\node[badge] at (8.08,-0.42) {6};
\node[hdr] at (8.28,-0.42) {template $\cfv_{\mathrm{tmpl}}$};
\foreach [count=\i from 0] \c in \Tfills \cell{\sxb+\i*\cw}{-0.95}{\c}
\seps{\sxb}{-0.95}{0, 1, 3, 7, 11}
\draw[arr] (7.45,-1.81) -- (7.70,-1.81) |- (7.90,-0.78);
% --- pass 5: depth schedule ---
\draw[arr] (7.45,-2.83) -- (7.90,-2.83);
\blockstrip{\sxb}{-3.00}{\Pfills}
\foreach \a/\b/\l in {7.95/8.63/0, 8.75/10.11/1, 10.23/11.59/2}{
  \draw[decorate, decoration={brace, mirror, amplitude=2pt}, black!55] (\a,-3.07) -- (\b,-3.07)
    node[note, midway, below=1.5pt] {$B_\l$};}
\draw[fold] (8.29,-2.96) -- (8.29,-0.99) node[note, pos=0.45, right=0.6mm] {fold chance};
\node[badge] at (8.08,-3.72) {5};
\node[hdr] at (8.28,-3.72) {depth schedule: stable sort by parent depth, $D{=}3$ blocks};
% --- pass 7: index emission ---
\def\sxc{12.45}
\node[badge] at (12.58,-0.42) {7};
\node[hdr] at (12.78,-0.42) {index emission per block (dark = $\bot$)};
\blockstrip{\sxc}{-0.95}{\Qone}
\blockstrip{\sxc}{-1.29}{\Qtwo}
\blockstrip{\sxc}{-1.63}{\Sfills}
\node[note, anchor=west] at (16.15,-0.78) {$q_1$};
\node[note, anchor=west] at (16.15,-1.12) {$q_2$};
\node[note, anchor=west] at (16.15,-1.46) {$q_{\mathrm{val}}$};
\draw[arr] (11.66,-2.83) -- (12.05,-2.83) |- (12.40,-1.29);
\draw[arr] (11.66,-2.83) -- (12.40,-2.83);
\foreach [count=\i from 0] \c in {teal!45, teal!45, teal!45, teal!45, violet!40, violet!40, violet!40, violet!40}
  \cell{\sxc+\i*\cw}{-3.00}{\c}
\node[hdr] at (12.45,-2.48) {flat over decision edges};
\foreach \a/\b/\l in {12.45/13.81/1, 13.81/15.17/2}{
  \draw[decorate, decoration={brace, mirror, amplitude=2pt}, black!55] (\a+0.02,-3.07) -- (\b-0.02,-3.07)
    node[note, midway, below=1.5pt] {player \l};}
% --- host to device: one copy ---
\draw[flow] (16.95,-4.18) -- (16.95,-5.40);
\node[note, anchor=east] at (16.80,-4.78) {copy to the device, once};

% ======================= device lane: passes 8-10 =======================
% the whole device lane is shifted down as one block so titles clear the frame
\begin{scope}[yshift=-0.45cm]
% --- pass 8: memory plan as bars proportional to element count, with the
%     pinned / mutable columns of Table memplan ---
\node[badge] at (0.35,-5.30) {8};
\node[hdr] at (0.55,-5.30) {memory plan};
\def\bx{1.95}      % bar origin; leave room for the row labels
\def\bs{0.085}     % cm per element on the example game
\def\px{4.85}      % pinned column
\def\mx{6.15}      % mutable column
\def\ypred{-8.75}  % shared prediction-hook baseline (uniform 0.60 row pitch)
\node[note, anchor=base] at (\px,-5.68) {pinned};
\node[note, anchor=base] at (\mx,-5.68) {mutable};
% rows: label / length in elements / fill / pinned / mutable
\foreach [count=\r from 0] \l/\n/\f/\p/\m in {
  {index arrays}/26/black!22/0/0,
  {$\cfv_{\mathrm{tmpl}}$}/11/white/1/1,
  {$\regret$}/8/orange!30/1/1,
  {$\bar{s}$ (f64)}/8/orange!15/1/1,
  {prediction}/8/orange!15/1/1,
  {$\strat_{\mathrm{ext}}$}/9/green!25/1/0,
  {$\reach_{1,2}$}/22/teal!28/1/0,
  {$\cfv$}/11/gray!15/1/0,
  {set totals}/4/gray!15/1/0,
  {counter $t$}/1/gray!15/1/1}{
  \pgfmathsetmacro{\y}{-5.95-\r*0.30}
  \node[note, anchor=east] at (\bx-0.10,\y) {\l};
  \draw[black!60, fill=\f] (\bx,\y-0.10) rectangle ++({\n*\bs},0.20);
  \ifnum\p=1 \fill[blue!60!black] (\px,\y) circle (1.4pt); \fi
  \ifnum\m=1 \draw[orange!75!black, line width=0.9pt] (\mx,\y) circle (1.6pt); \fi
}
% index arrays are ~9E elements: draw the bar broken and annotate
\fill[white] (3.35,-6.13) -- (3.42,-5.77) -- (3.52,-5.77) -- (3.45,-6.13) -- cycle;
\draw[black!60] (3.35,-6.13) -- (3.42,-5.77); \draw[black!60] (3.45,-6.13) -- (3.52,-5.77);
\node[note, anchor=west] at (4.02,-5.95) {${\approx}9\numedges$};

% --- pass 9: rule lowering, hooks bound into the iteration body ---
\node[badge] at (6.95,-5.30) {9};
\node[hdr] at (7.15,-5.30) {rule lowering};
\node[draw=black!30, rounded corners=4pt, fill=black!2, minimum width=20mm,
      minimum height=33mm, anchor=north] (body) at (8.10,-5.95) {};
\node[hook] at (8.10,-6.35) {discount};
\node[note] at (8.10,-6.95) {forward};
\node[note] at (8.10,-7.55) {backward};
\node[hook] at (8.10,-8.15) {after};
\node[hook] at (8.10,\ypred) {prediction};
\node[note, anchor=north] at ($(body.south)+(0,-0.05)$) {iteration body};
% rule columns: which hooks each shipped rule binds (the four rules of Table 19)
\def\rx{9.75}
\foreach [count=\j from 0] \r in {CFR, CFR$^{+}$, PCFR$^{+}$, DCFR}
  \node[note, rotate=60, anchor=west, inner sep=0pt] at ({\rx+\j*0.55},-6.10) {\r};
\foreach \y in {-6.35,-8.15,\ypred} \draw[black!25] (9.15,\y) -- (11.65,\y);
\foreach \j in {0,...,3} \foreach \y in {-6.35,-8.15,\ypred}
  \draw[black!35, fill=white] ({\rx+\j*0.55},\y) circle (1.5pt);
% (rule, hook row) filled where bound: discount y=-6.35, after -8.15, prediction \ypred
\foreach \j/\y in {1/-8.15, 2/-6.35, 2/-8.15, 2/\ypred, 3/-6.35}
  \fill[orange!85!black] ({\rx+\j*0.55},\y) circle (1.5pt);

% --- pass 10: runtime; warm up, record once, replay every iteration ---
\node[badge] at (12.35,-5.30) {10};
\node[hdr] at (12.60,-5.30) {runtime};
\def\wz{13.35}
\node[note, anchor=east] at (\wz-0.08,-6.05) {warm-up};
\draw[black!35] (\wz,-6.05) -- (16.05,-6.05);
\foreach \i in {0,1,2}{
  \foreach \k in {0,...,5}{
    \draw[tick] ({\wz+0.10+\i*0.90+\k*0.13},-5.92) -- ++(0,-0.26);}}
\node[draw=orange!80!black, fill=orange!25, rounded corners=1.5pt,
      minimum height=3.4mm, inner sep=1.6pt, anchor=west] (rec) at (16.15,-6.05) {record};
\node[note, anchor=east] at (\wz-0.08,-7.05) {replay};
\draw[black!35] (\wz,-7.05) -- (16.35,-7.05);
\foreach \i in {0,...,3}
  \draw[gtick] ({\wz+0.20+\i*0.90},-6.86) -- ++(0,-0.38);
\draw[flow, line width=0.9pt] (16.35,-7.05) -- ++(0.30,0);
\node[note, anchor=west] at (16.70,-7.05) {$\times T$};
% between solves: only the mutable buffers are rewritten; the graph is kept
\draw[ret] (16.35,-7.65) -- (16.35,-9.95) -- (\mx,-9.95) -- (\mx,-8.90);
\node[note, anchor=south east, text=orange!60!black, align=right] at (16.25,-9.90)
  {between solves:\\ rewrite mutable buffers only;\\ the captured graph stays valid};
\end{scope}

% ======================= lane frames =======================
\begin{scope}[on background layer]
\node[draw=blue!45!black!30, rounded corners=3pt, inner sep=1.2mm, fill=blue!3,
      fit={(0.10,-0.25) (17.40,-3.98)}] (hostlane) {};
\node[draw=orange!70!black!30, rounded corners=3pt, inner sep=1.2mm, fill=orange!4,
      fit={(0.10,-5.55) (17.40,-10.55)}] (devlane) {};
\end{scope}
\node[lanelabel, anchor=south west] at (hostlane.north west)
  {Build on the host, once per game (\cref{alg:compile}, lines 1--7)};
\node[lanelabel, anchor=south west] at (devlane.north west)
  {Device: allocate once, replay every iteration (lines 8--10)};
\end{tikzpicture}}
\caption{The compiler passes of \cref{alg:compile} on the example game of
\cref{fig:archpipeline}; badges are algorithm lines. The host lane flattens
the tree into slot, node, and edge strips, sorts the edges into $D$ depth
blocks, folds chance into the values template, and emits the per-lane index
arrays, where dark cells point at the sentinel $\bot$. The device lane sizes
every buffer of \cref{tab:memplan} once (bars scale with element count; dots
mark pinned addresses, rings the inputs rewritten between solves), binds the
hooks each update rule uses (dots: bound), and after three eager iterations
records one graph that is replayed thereafter; between solves only the ringed
buffers are rewritten.}
\Description{Two stacked lanes. The top lane shows a small game tree passing a
diamond-shaped validation gate into three flat strips of colored cells for
slots, nodes, and edges, then an arrow into the same edges regrouped into three
depth blocks with braces, a dashed arrow folding chance into a template strip
whose terminal cells are dark, and finally three index strips per block where
opponent and chance cells are dark. A thick arrow drops into the bottom lane
labeled copy once. The bottom lane has a bar chart of buffer sizes with a
pinned column of blue dots and a mutable column of orange rings, a rounded
capsule listing the iteration body with three dashed hook slots and a dot
matrix of four update rules showing which hooks each binds, and two timelines:
a warm-up row with dense ticks, a record box, and a replay row with one thick
tick per iteration. A dashed orange return path leads from the replay row back
to the mutable column.}
\label{fig:compiler}
\end{figure*}


\paragraph{Planning scope.}
The memory plan is an estimate: it sizes the scheduled buffers from the
game's counts and reports whether they fit the configured budget, and the
solver reads that verdict as a planning indicator and builds regardless. The regret and average-strategy vectors are allocated before the
schedule is planned, so the passes above describe the logical order of the
build.

\paragraph{Input representation.}
The input is a tabular game specification: a map from node id to (type, acting
player, information-set key, action list with child ids, terminal payoff), a map
from information-set key to (player, action names), and a map from node id to the
chance probability of reaching it from its parent. Native games and the
OpenSpiel and file-based converters all produce this one structure, so every
pass below it is shared; a converter is validated by the structural and
behavioral checks of \cref{app:protocol}.

\paragraph{Mutable and immutable fields.}
\Cref{tab:memplan} separates what the graph pins from what a caller may change.
Every index array is immutable after emission. Two inputs are mutable in place
without recompilation: the terminal payoffs and the root chance distribution,
both of which enter the iteration only through the values template
(\cref{eq:template}), which \texttt{set\_payoffs} and \texttt{set\_root\_chance}
rewrite in place from the retained per-node root owner and below-root chance
product. The solver state is reset in place by \texttt{reset()}, which zeroes
the two accumulators and the prediction buffer and resets the iteration counter;
because no buffer is reallocated, the captured graph stays valid and the next
call replays from iteration one. \Cref{tab:lifecycle} reports the cost of each
of these operations.

\begin{table*}[!htb]
\centering
\caption{Memory plan of the compiled solver. Sizes are element counts; the
dtype is the configured compute dtype except where noted. \emph{Pinned} marks
buffers whose addresses the captured graph records; \emph{Mutable} marks the
inputs a caller may rewrite in place between solves.}
\label{tab:memplan}
{
\setlength{\tabcolsep}{4pt}
\begin{tabular}{llccl}
\toprule
Buffer & Size & Pinned & Mutable & Written \\
\midrule
node arrays (type, owner, set, offsets) & $\numnodes$, int & -- & -- & build \\
edge index arrays (all of step 7) & $\approx 7\numedges + 2\numedges$, int64 & -- & -- & build \\
chance reach $\pi_c$ & $\numnodes$ & -- & via root & build \\
values template $\cfv_{\mathrm{tmpl}}$ & $\numnodes$ & yes & yes & build, updates \\
regrets $\regret$ & $\numslots$ & yes & reset & every iteration \\
strategy sum $\bar{s}$ (float64) & $\numslots$ & yes & reset & every iteration \\
prediction buffer (PCFR$^{+}$) & $\numslots$ & yes & reset & every iteration \\
strategy $\strat_{\mathrm{ext}}$ (+ sentinel) & $\numslots + 1$ & yes & -- & every iteration \\
dual-lane reach & $2\numnodes$ & yes & -- & every iteration \\
values $\cfv$ & $\numnodes$ & yes & -- & every iteration \\
set totals & $\numinfosets$ & yes & -- & every iteration \\
iteration counter $t$ (device) & 1, float32 & yes & reset & every replay \\
CUDA-graph pool & driver-managed & -- & -- & capture \\
\bottomrule
\end{tabular}}
\end{table*}

\paragraph{Update rules share one schedule.}
\Cref{tab:rules} counts the Aten operations each shipped rule issues per
iteration. The rules differ only in the hooks of step 9 (the dot matrix of
\cref{fig:compiler}), which act on the slot vectors: the delta against CFR$^{+}$ is a constant per rule (from $-1$ for vanilla
CFR, which drops the clamp, to $+29$ for DCFR's discount schedule) and
does not change between Leduc and the 434k-state turn subgame, while the
tree-dependent part of the schedule, the $8D$ operations of the two passes, is
shared verbatim. Every rule's body captures into a CUDA graph, because the hooks
read the iteration index from the device-side counter. The convergence of the rules is reported in \cref{app:extra}.

\input{figures/TABLE_T16_rules.tex}
